
\documentclass{article}
\usepackage{colortbl}
\usepackage{makecell}
\usepackage{multirow}
\usepackage{supertabular}

\begin{document}

\newcounter{utterance}

\centering \large Interaction Transcript for game `hot\_air\_balloon', experiment `air\_balloon\_survival\_en\_negotiation\_easy', episode 1 with Qwen3{-}32B{-}t0.0.
\vspace{24pt}

{ \footnotesize  \setcounter{utterance}{1}
\setlength{\tabcolsep}{0pt}
\begin{supertabular}{c@{$\;$}|p{.15\linewidth}@{}p{.15\linewidth}p{.15\linewidth}p{.15\linewidth}p{.15\linewidth}p{.15\linewidth}}
    \# & $\;$A & \multicolumn{4}{c}{Game Master} & $\;\:$B\\
    \hline

    \theutterance \stepcounter{utterance}  
    & & \multicolumn{4}{p{0.6\linewidth}}{
        \cellcolor[rgb]{0.9,0.9,0.9}{
            \makecell[{{p{\linewidth}}}]{
                \texttt{\tiny{[P1$\langle$GM]}}
                \texttt{You are participating in a collaborative negotiation game.} \\
\\ 
\texttt{Together with another participant, you must agree on a single set of items that will be kept. Each of you has your own view of how much each item matters to you (importance). You do not know how the other participant values the items. Additionally, you are given the effort each item demands.} \\
\texttt{You may only agree on a set if the total effort of the selected items does not exceed a shared limit:} \\
\\ 
\texttt{LIMIT = 4564} \\
\\ 
\texttt{Here are the individual item effort values:} \\
\\ 
\texttt{Item effort = \{"C08": 299, "C55": 241, "A43": 889, "A67": 967, "B42": 619, "A62": 799, "B19": 978, "C43": 733, "C73": 909, "C10": 501, "B36": 139, "A18": 594, "C61": 565, "C67": 789, "C42": 107\}} \\
\\ 
\texttt{Here is your personal view on the importance of each item:} \\
\\ 
\texttt{Item importance values = \{"C08": 138, "C55": 583, "A43": 868, "A67": 822, "B42": 783, "A62": 65, "B19": 262, "C43": 121, "C73": 508, "C10": 780, "B36": 461, "A18": 484, "C61": 668, "C67": 389, "C42": 808\}} \\
\\ 
\texttt{Goal:} \\
\\ 
\texttt{Your goal is to negotiate a shared set of items that benefits you as much as possible (i.e., maximizes total importance to YOU), while staying within the LIMIT. You are not required to make a PROPOSAL in every message {-} you can simply negotiate as well. All tactics are allowed!} \\
\\ 
\texttt{Interaction Protocol:} \\
\\ 
\texttt{You may only use the following structured formats in a message:} \\
\\ 
\texttt{PROPOSAL: \{'A', 'B', 'C', …\}} \\
\texttt{Propose keeping exactly those items.} \\
\texttt{REFUSE: \{'A', 'B', 'C', …\}} \\
\texttt{Explicitly reject opponent's proposal.} \\
\texttt{ARGUMENT: \{'...'\}} \\
\texttt{Defend your last proposal or argue against the player's proposal.} \\
\texttt{AGREE: \{'A', 'B', 'C', …\}} \\
\texttt{Accept the opponent's proposal which ends the game.} \\
\texttt{STRATEGIC REASONING: \{'...'\}} \\
\texttt{	Describe your strategic reasoning or anticipation explaining your choice of action. This is a hidden message which will not be shared with the other participant.} \\
\\ 
\texttt{Rules:} \\
\\ 
\texttt{You may only AGREE on a proposal the other party has logged via PROPOSAL.} \\
\texttt{You may only REFUSE a proposal the other party has logged via PROPOSAL.} \\
\texttt{Total effort of any PROPOSAL or AGREE set must be ≤ LIMIT.} \\
\texttt{Do NOT reveal your hidden importance scores.} \\
\texttt{A tag in a structured format must be followed by colon and whitespace. The argument must be a python set containing 0 or more strings.} \\
\texttt{So, it must be of the form TAG: \{...\}} \\
\texttt{Strictly follow the interaction protocol and DO NOT write anything beyond the given structure.} \\
\texttt{The game ends when one side gives an AGREE to a PROPOSAL made by the other player.} \\
\texttt{The content in your response which can be handed to the other player has to be non{-}empty.} \\
\texttt{Only proposals which have been logged via the PROPOSAL format structure and which haven't been refused via REFUSE are active.} \\
\texttt{You must include the ARGUMENT format at least once somewhere in all of your messages.} \\
\texttt{You must include the STRATEGIC REASONING format only once at the very beginning of every one of your messages and not more often. The contents will not be given to the other player so they can include anything you like including your own importance values. Here you should reason step by step to come up with you next move.} \\
            }
        }
    }
    & \\ \\

    \theutterance \stepcounter{utterance}  
    & \multicolumn{4}{p{0.6\linewidth}}{
        \cellcolor[rgb]{0.9,0.9,0.9}{
            \makecell[{{p{\linewidth}}}]{
                \texttt{\tiny{[P1$\rangle$GM]}}
                \texttt{<think>} \\
\texttt{Okay, let's see. I need to figure out the best set of items to propose that maximizes my total importance without exceeding the effort limit of 4564. First, I should calculate the total effort of all items to see if the sum is over the limit. Let me add them up.} \\
\\ 
\texttt{Looking at the item efforts: C08 is 299, C55 is 241, A43 is 889, A67 is 967, B42 is 619, A62 is 799, B19 is 978, C43 is 733, C73 is 909, C10 is 501, B36 is 139, A18 is 594, C61 is 565, C67 is 789, C42 is 107. Let me sum these up step by step.} \\
\\ 
\texttt{Starting with 0:} \\
\texttt{+299 = 299} \\
\texttt{+241 = 540} \\
\texttt{+889 = 1429} \\
\texttt{+967 = 2396} \\
\texttt{+619 = 3015} \\
\texttt{+799 = 3814} \\
\texttt{+978 = 4792. Oh wait, that's already over 4564. So the total effort of all items is way over the limit. So we can't take all. Need to find a subset.} \\
\\ 
\texttt{My goal is to maximize my own importance. Let me list the items with their importance and effort:} \\
\\ 
\texttt{C08: 138 (299)} \\
\texttt{C55: 583 (241)} \\
\texttt{A43: 868 (889)} \\
\texttt{A67: 822 (967)} \\
\texttt{B42: 783 (619)} \\
\texttt{A62: 65 (799)} \\
\texttt{B19: 262 (978)} \\
\texttt{C43: 121 (733)} \\
\texttt{C73: 508 (909)} \\
\texttt{C10: 780 (501)} \\
\texttt{B36: 461 (139)} \\
\texttt{A18:} \\
            }
        }
    }
    & & \\ \\

    \theutterance \stepcounter{utterance}  
    & & \multicolumn{4}{p{0.6\linewidth}}{
        \cellcolor[rgb]{0.9,0.9,0.9}{
            \makecell[{{p{\linewidth}}}]{
                \texttt{\tiny{[P1$\langle$GM]}}
                \texttt{Your response did not start with the proper strategic reasoning tag at the very beginning of your response. The very first structured format must be of the form STRATEGIC REASONING: \{...\}. Try again.} \\
            }
        }
    }
    & \\ \\

    \theutterance \stepcounter{utterance}  
    & \multicolumn{4}{p{0.6\linewidth}}{
        \cellcolor[rgb]{0.9,0.9,0.9}{
            \makecell[{{p{\linewidth}}}]{
                \texttt{\tiny{[P1$\rangle$GM]}}
                \texttt{<think>} \\
\texttt{Okay, let's start over. I need to make sure I follow the correct format. First, I should calculate the total effort of all items to see if it's over the limit. The LIMIT is 4564. Let me add up all the item efforts:} \\
\\ 
\texttt{C08: 299} \\
\texttt{C55: 241} \\
\texttt{A43: 889} \\
\texttt{A67: 967} \\
\texttt{B42: 619} \\
\texttt{A62: 799} \\
\texttt{B19: 978} \\
\texttt{C43: 733} \\
\texttt{C73: 909} \\
\texttt{C10: 501} \\
\texttt{B36: 139} \\
\texttt{A18: 594} \\
\texttt{C61: 565} \\
\texttt{C67: 789} \\
\texttt{C42: 107} \\
\\ 
\texttt{Adding them up step by step:} \\
\\ 
\texttt{299 + 241 = 540} \\
\texttt{540 + 889 = 1429} \\
\texttt{1429 + 967 = 2396} \\
\texttt{2396 + 619 = 3015} \\
\texttt{3015 + 799 = 3814} \\
\texttt{3814 + 978 = 4792 (already over the limit here)} \\
\\ 
\texttt{So the total effort is 4792, which is over 4564. Therefore, we can't take all items. I need to find a subset that maximizes my importance without exceeding the limit.} \\
\\ 
\texttt{My importance values are:} \\
\\ 
\texttt{C08: 138} \\
\texttt{C55: 583} \\
\texttt{A43: 868} \\
\texttt{A67: 822} \\
\texttt{B42: 783} \\
\texttt{A62: 65} \\
\texttt{B19: 262} \\
\texttt{C43: 121} \\
\texttt{C73: 508} \\
\texttt{C10: 780} \\
\texttt{B36: 461} \\
\texttt{A18: 484} \\
\texttt{C61: 668} \\
\texttt{C67: 389} \\
\texttt{C42: 808} \\
\\ 
\texttt{I should prioritize items with high importance per unit of effort. Let's calculate the ratio} \\
            }
        }
    }
    & & \\ \\

    \theutterance \stepcounter{utterance}  
    & & \multicolumn{4}{p{0.6\linewidth}}{
        \cellcolor[rgb]{0.9,0.9,0.9}{
            \makecell[{{p{\linewidth}}}]{
                \texttt{\tiny{[P1$\langle$GM]}}
                \texttt{Your response did not start with the proper strategic reasoning tag at the very beginning of your response. The very first structured format must be of the form STRATEGIC REASONING: \{...\}. Try again.} \\
            }
        }
    }
    & \\ \\

    \theutterance \stepcounter{utterance}  
    & \multicolumn{4}{p{0.6\linewidth}}{
        \cellcolor[rgb]{0.9,0.9,0.9}{
            \makecell[{{p{\linewidth}}}]{
                \texttt{\tiny{[P1$\rangle$GM]}}
                \texttt{<think>} \\
\texttt{Okay, I need to make sure I start with the STRATEGIC REASONING tag. Let me outline my plan step by step.} \\
\\ 
\texttt{First, I'll calculate the total effort of all items to see if it's over the limit. The LIMIT is 4564. Adding up all the item efforts:} \\
\\ 
\texttt{C08: 299} \\
\texttt{C55: 241} \\
\texttt{A43: 889} \\
\texttt{A67: 967} \\
\texttt{B42: 619} \\
\texttt{A62: 799} \\
\texttt{B19: 978} \\
\texttt{C43: 733} \\
\texttt{C73: 909} \\
\texttt{C10: 501} \\
\texttt{B36: 139} \\
\texttt{A18: 594} \\
\texttt{C61: 565} \\
\texttt{C67: 789} \\
\texttt{C42: 107} \\
\\ 
\texttt{Total effort: 299 + 241 = 540; 540 + 889 = 1429; 1429 + 967 = 2396; 2396 + 619 = 3015; 3015 + 799 = 3814; 3814 + 978 = 4792. That's already over 4564. So we can't take all items.} \\
\\ 
\texttt{Next, I need to maximize my importance. My importance values are:} \\
\\ 
\texttt{C08: 138} \\
\texttt{C55: 583} \\
\texttt{A43: 868} \\
\texttt{A67: 822} \\
\texttt{B42: 783} \\
\texttt{A62: 65} \\
\texttt{B19: 262} \\
\texttt{C43: 121} \\
\texttt{C73: 508} \\
\texttt{C10: 780} \\
\texttt{B36: 461} \\
\texttt{A18: 484} \\
\texttt{C61: 668} \\
\texttt{C67: 389} \\
\texttt{C42: 808} \\
\\ 
\texttt{I should prioritize items with high importance per unit of effort. Let's calculate the ratio for each item:} \\
\\ 
\texttt{C08: 138/2} \\
            }
        }
    }
    & & \\ \\

    \theutterance \stepcounter{utterance}  
    & & & \multicolumn{2}{p{0.3\linewidth}}{
        \cellcolor[rgb]{0.9,0.9,0.9}{
            \makecell[{{p{\linewidth}}}]{
                \texttt{\tiny{[GM$|$GM]}}
                \texttt{abort} \\
            }
        }
    }
    & & \\ \\

    \theutterance \stepcounter{utterance}  
    & & & \multicolumn{2}{p{0.3\linewidth}}{
        \cellcolor[rgb]{0.9,0.9,0.9}{
            \makecell[{{p{\linewidth}}}]{
                \texttt{\tiny{[GM$|$GM]}}
                \texttt{end game} \\
            }
        }
    }
    & & \\ \\

\end{supertabular}
}

\end{document}
